% ══════════════════════════════════════════════════════════════════════════════
%  ObjectIR v2 — Language Specification
% ══════════════════════════════════════════════════════════════════════════════
\documentclass[11pt,a4paper]{article}

% ── Packages ──────────────────────────────────────────────────────────────────
\usepackage[T1]{fontenc}
\usepackage[utf8]{inputenc}
\usepackage{geometry}
\usepackage{hyperref}
\usepackage{xcolor}
\usepackage{listings}
\usepackage{longtable}
\usepackage{booktabs}
\usepackage{array}
\usepackage{enumitem}
\usepackage{titlesec}
\usepackage{fancyhdr}
\usepackage{parskip}
\usepackage{microtype}
\usepackage{lmodern}
\usepackage{tabularx}
\usepackage{mdframed}
\usepackage{amsmath}

\geometry{margin=2.4cm}

% ── Colors ────────────────────────────────────────────────────────────────────
\definecolor{codebg}{HTML}{F5F5F5}
\definecolor{codeborder}{HTML}{CCCCCC}
\definecolor{kw}{HTML}{0000CC}
\definecolor{str}{HTML}{008800}
\definecolor{comment}{HTML}{888888}
\definecolor{typecolor}{HTML}{7B0000}
\definecolor{linkcolor}{HTML}{0066CC}
\definecolor{sectionrule}{HTML}{3A6EA5}
\definecolor{notebg}{HTML}{FFF8DC}
\definecolor{noteborder}{HTML}{DAA520}
\definecolor{warnbg}{HTML}{FFF0F0}
\definecolor{warnborder}{HTML}{CC4444}
\definecolor{infobg}{HTML}{EEF4FF}
\definecolor{infoborder}{HTML}{5577CC}

% ── Hyperref ──────────────────────────────────────────────────────────────────
\hypersetup{
  colorlinks=true,
  linkcolor=linkcolor,
  urlcolor=linkcolor,
  citecolor=linkcolor,
  pdftitle={ObjectIR V3 --- Language Specification},
  pdfauthor={Finite RnD},
}

% ── Code listing styles ───────────────────────────────────────────────────────
\lstdefinelanguage{TextIR}{
  morekeywords={module,version,class,interface,struct,enum,field,method,
                constructor,local,public,private,protected,internal,static,
                virtual,override,abstract,sealed,implements,else,catch,finally,
                true,false,null,ldc,i4,.,i8,ldstr,ldnull,stloc,ldloc,ldfld,stfld,
                ldelem,stelem,add,sub,mull,div},
  morestring=[b]",
  morecomment=[l]{//},
  sensitive=true
}

\lstdefinestyle{textir}{
  language=TextIR,
  backgroundcolor=\color{codebg},
  frame=single,
  rulecolor=\color{codeborder},
  basicstyle=\ttfamily\small,
  keywordstyle=\bfseries\color{kw},
  stringstyle=\color{str},
  commentstyle=\itshape\color{comment},
  breaklines=true,
  showstringspaces=false,
  tabsize=4,
  numbers=left,
  numberstyle=\tiny\color{comment},
  numbersep=6pt,
  xleftmargin=1.8em,
}

\lstdefinestyle{pseudo}{
  backgroundcolor=\color{codebg},
  frame=single,
  rulecolor=\color{codeborder},
  basicstyle=\ttfamily\small,
  breaklines=true,
  showstringspaces=false,
  tabsize=4,
  numbers=none,
}

\lstdefinestyle{json}{
  backgroundcolor=\color{codebg},
  frame=single,
  rulecolor=\color{codeborder},
  basicstyle=\ttfamily\small,
  stringstyle=\color{str},
  breaklines=true,
  showstringspaces=false,
  tabsize=2,
  numbers=none,
}

\newcommand{\code}[1]{\texttt{\small #1}}
\newcommand{\typename}[1]{\textcolor{typecolor}{\texttt{\small #1}}}
\newcommand{\kword}[1]{\textbf{\texttt{\small #1}}}
\newcommand{\opcode}[1]{\textbf{\texttt{#1}}}
\newcommand{\nterm}[1]{\textit{$\langle$#1$\rangle$}}
\newcommand{\term}[1]{\texttt{#1}}

% ── Callout boxes ─────────────────────────────────────────────────────────────
\newmdenv[
  backgroundcolor=notebg, linecolor=noteborder, linewidth=1pt,
  innerleftmargin=8pt, innerrightmargin=8pt,
  innertopmargin=6pt, innerbottommargin=6pt,
  skipabove=6pt, skipbelow=6pt,
]{notebox}

\newmdenv[
  backgroundcolor=warnbg, linecolor=warnborder, linewidth=1pt,
  innerleftmargin=8pt, innerrightmargin=8pt,
  innertopmargin=6pt, innerbottommargin=6pt,
  skipabove=6pt, skipbelow=6pt,
]{warnbox}

\newmdenv[
  backgroundcolor=infobg, linecolor=infoborder, linewidth=1pt,
  innerleftmargin=8pt, innerrightmargin=8pt,
  innertopmargin=6pt, innerbottommargin=6pt,
  skipabove=6pt, skipbelow=6pt,
]{infobox}

% ── Headers / footers ─────────────────────────────────────────────────────────
\pagestyle{fancy}
\fancyhf{}
\renewcommand{\headrulewidth}{0.4pt}
\fancyhead[L]{\small\textit{ObjectIR V3 --- Language Specification}}
\fancyhead[R]{\small\textit{ObjectIR Project}}
\fancyfoot[C]{\thepage}

% ── Section formatting ────────────────────────────────────────────────────────
\titleformat{\section}{\large\bfseries\color{sectionrule}}{}{0em}{}[{\color{sectionrule}\titlerule[0.8pt]}]
\titleformat{\subsection}{\normalsize\bfseries}{}{0em}{}
\titleformat{\subsubsection}{\normalsize\itshape\bfseries}{}{0em}{}

% ══════════════════════════════════════════════════════════════════════════════
\begin{document}

\begin{titlepage}
  \centering
  \vspace*{3cm}
  {\Huge\bfseries ObjectIR}\\[0.5em]
  {\large\bfseries Version 3}\\[0.8em]
  {\Large Language Specification}\\[2.5em]
  \rule{0.55\linewidth}{0.4pt}\\[1.2em]
  {\normalsize
    A stack-based intermediate representation\\[0.3em]
    and virtual machine definition
  }\\[2.5em]
  \begin{tabular}{ll}
    \textbf{Source format:}   & TextIR (human-readable), JSON, FOB/IR v3 (binary) \\[0.2em]
    \textbf{Execution model:} & Stack machine, one evaluation stack per call frame \\[0.2em]
    \textbf{Type model:}      & Untyped stack; lazy coercion at instruction boundaries \\[0.2em]
    \textbf{Revision:}        & \today \\
  \end{tabular}
  \vfill
  {\small This document is the normative specification of the ObjectIR language:
  its source format, module structure, type system, and instruction set.
  Implementation-specific notes are confined to appendices.}
\end{titlepage}

\tableofcontents
\newpage

% ══════════════════════════════════════════════════════════════════════════════
\section{Introduction}
% ══════════════════════════════════════════════════════════════════════════════

\textbf{ObjectIR} is a typed, object-oriented intermediate representation (IR)
designed to be simple to parse, simple to execute, and straightforward to
target from higher-level languages or toolchains.

\begin{notebox}
	All of the examples shown in this document describe the TextIR variant of ObjectIR. see section~\ref{sec:TextIRFeatures} for more information about what TextIR is.
\end{notebox}


A program is expressed as one or more \emph{modules}, each containing class
and interface definitions whose methods are composed of explicit stack-machine
instructions. The IR sits at the same conceptual level as JVM bytecode or
CIL, but its surface format (\emph{TextIR}) is human-readable and
hand-writable.

\begin{notebox}
people working should use the ~\ref{sec:fob} Fob format for proper use cases and TextIR should only be used when 
\end{notebox}
\subsection{Design Goals}

\begin{itemize}[leftmargin=1.8em]
  \item \textbf{Human-readable source.} TextIR is the canonical format: a
    whitespace-significant, keyword-driven text form that lowers to the same
    module representation as the JSON and binary forms.
  \item \textbf{Multiple input formats.} A module may be delivered as TextIR
    source, a JSON object graph, or a compact FOB/IR v3 binary. All three
    forms are semantically equivalent.
  \item \textbf{Embeddable.} A conforming runtime must allow host code to
    load a module, invoke individual methods, and register external (native)
    type and method bindings without touching a parser.
  \item \textbf{Lazy coercion.} Values on the evaluation stack carry no
    static type tag. Numeric widening and other conversions are performed
    at instruction boundaries according to well-defined rules.
  \item \textbf{Structured control flow.} Loops, conditionals, and exception
    handlers are first-class instructions, not jump targets.
\end{itemize}

\subsection{Document Conventions}

\begin{itemize}[leftmargin=1.8em]
  \item \textbf{Normative} text uses \textbf{shall} / \textbf{must}.
    \textbf{Informative} text uses \textit{should} / \textit{may}.
  \item Grammar non-terminals are written \nterm{like-this}; terminals are
    written \term{like-this}.
  \item Instruction mnemonics are typeset \opcode{like-this}.
  \item Type names are typeset \typename{like-this}.
  \item this document describes information for "conforming" runtime's.
  not everything is covered in this document. for more information, please refer to the reference runtime for wishing to use ObjectIR with a runtime instead of transpilation.
\end{itemize}

% ══════════════════════════════════════════════════════════════════════════════
\section{Module Structure}
\label{sec:module}
% ══════════════════════════════════════════════════════════════════════════════

A \emph{module} is the top-level compilation and loading unit of ObjectIR.

\subsection{Module Contents}

A well-formed module contains:

\begin{itemize}[leftmargin=1.8em]
  \item A \textbf{name}: a non-empty identifier string.
  \item A \textbf{version}: a three-part numeric tuple
    (\textit{major}.\textit{minor}.\textit{patch}), where \textit{patch} is
    optional and defaults to \texttt{0}.
  \item An ordered list of \textbf{type definitions}: classes, interfaces,
    structs, and enums.
  \item A (reserved) \textbf{free-function table}, intended for future
    module-level functions outside any type. 
\end{itemize}

\subsection{Module Identity}

Two modules are considered the \emph{same module} if and only if their names
are identical. Version numbers are informational; the runtime does not perform
version enforcement unless the host explicitly requires it.

\subsection{Active Module}

A runtime instance holds \textbf{exactly one} active module at a time.
Loading a new module replaces the previous one. A \emph{full reset} also
clears all static field state; a \emph{module-only reload} replaces the type
table while optionally preserving heap state.

% ══════════════════════════════════════════════════════════════════════════════
\section{Source Formats}
\label{sec:formats}
% ══════════════════════════════════════════════════════════════════════════════

ObjectIR defines three equivalent surface representations of a module. All
three must express identical semantics; a conforming tool or runtime may
accept any subset.

\begin{longtable}{lp{10cm}}
  \toprule
  \textbf{Format} & \textbf{Description} \\
  \midrule
  \endhead
  \textbf{TextIR}    &
    Human-readable source text (see Section~\ref{sec:textir}).
    The canonical authoring format. A conforming parser produces a module
    representation equivalent to the JSON form. \\[0.4em]
  \textbf{JSON}      &
    A JSON object whose structure mirrors the module schema
    (Section~\ref{sec:schema}). Property names are camelCase.
    This is the recommended interchange format between tools. \\[0.4em]
  \textbf{FOB/IR v3} &
    A compact binary encoding of the module schema
    (Section~\ref{sec:fob}). Intended for distribution and fast loading. \\
  \bottomrule
\end{longtable}

\begin{notebox}
\textbf{Auto-detection hint:} A string whose first non-whitespace character
is \texttt{\{} or \texttt{[} should be treated as JSON; all other strings
should be treated as TextIR. Implementations may expose this heuristic as
an \texttt{Auto} format option.
\end{notebox}

% ══════════════════════════════════════════════════════════════════════════════
\section{TextIR Syntax}
\label{sec:textir}
% ══════════════════════════════════════════════════════════════════════════════

TextIR is a line-oriented, brace-delimited language. Whitespace (outside
string literals) is insignificant except as a token separator. Comments
begin with \term{//} and extend to the end of the line.

\subsection{Module Header}

Every TextIR file must begin with a module header:

\begin{lstlisting}[style=textir]
module <name> version <major>.<minor>[.<patch>]
\end{lstlisting}

\subsection{Type Declarations}

\begin{center}
\begin{tabular}{lll}
  \toprule
  \textbf{Keyword}   & \textbf{Kind}      & \textbf{Notes} \\
  \midrule
  \term{class}     & Reference type   & Instances are heap-allocated objects. \\
  \term{interface} & Contract type    & Declares virtual methods; no fields. \\
  \term{struct}    & Value type       & Semantics are host-defined. \\
  \term{enum}      & Enumeration type & Members are named integer constants. \\
  \bottomrule
\end{tabular}
\end{center}

Any combination of the following \emph{access modifiers} may precede the
type keyword: \term{public}, \term{private}, \term{protected},
\term{internal}.

The following \emph{type modifiers} may also precede the keyword:
\term{abstract}, \term{sealed}.

Inheritance and interface implementation are expressed with \term{:}:

\begin{lstlisting}[style=textir]
class Animal {
    field name: string
    virtual method Speak() -> void { ret }
}

class Dog : Animal {
    override method Speak() -> void implements Animal.Speak {
        ldstr "Woof!"
        call System.Console.WriteLine(string) -> void
        ret
    }
}
\end{lstlisting}

\subsection{Member Declarations}

\begin{lstlisting}[style=textir]
// Field
[static] field <name>: <type>

// Method
[static] [virtual] [override] [abstract]
method <name>(<params>) -> <rettype> [implements <Interface>.<MethodName>]
{
    [local <name>: <type>]*
    <instructions>
}

// Constructor
constructor(<params>) {
    [local <name>: <type>]*
    <instructions>
}
\end{lstlisting}

Parameter lists and local variable declarations share the syntax
\code{name: type}, comma-separated.


\subsection{Attributes (Annotations)}

Custom attributes may be applied to types, fields, and methods by placing
\term{@AttributeName} (with an optional positional argument list) immediately
before the declaration:

\begin{lstlisting}[style=textir]
@Serializable
class PlayerData {
    @NonSerialized
    field cachedHash: int32

    @Deprecated("Use SaveAsync instead")
    method Save() -> void { ret }
}
\end{lstlisting}

A runtime must attach attribute metadata to the corresponding declaration;
how it is stored and queried is implementation-defined.

\subsection{String Escape Sequences}

Within string literals the following escape sequences are defined:

\begin{center}
\begin{tabular}{ll}
  \toprule
  \textbf{Sequence} & \textbf{Meaning} \\
  \midrule
  \texttt{\textbackslash n}              & Line feed (U+000A) \\
  \texttt{\textbackslash t}              & Horizontal tab (U+0009) \\
  \texttt{\textbackslash r}              & Carriage return (U+000D) \\
  \texttt{\textbackslash\textbackslash}  & Backslash (U+005C) \\
  \texttt{\textbackslash"}              & Double quote (U+0022) \\
  \bottomrule
\end{tabular}
\end{center}

% ══════════════════════════════════════════════════════════════════════════════
\section{Type System}
\label{sec:types}
% ══════════════════════════════════════════════════════════════════════════════

\subsection{Primitive Types}

ObjectIR defines a fixed set of primitive types. Type names in operands are
\textbf{case-insensitive}: the runtime normalises them to lower-case before
lookup, so \texttt{System.String}, \texttt{system.string}, and \texttt{string}
are all equivalent.

\begin{longtable}{lll}
  \toprule
  \textbf{Primary IR name(s)} & \textbf{Alias(es)} & \textbf{Width / Semantics} \\
  \midrule
  \endhead
  \texttt{void}                     & \texttt{system.void}     & No value; return type only. \\
  \texttt{bool}                     & \texttt{system.boolean}  & Boolean true/false. \\
  \texttt{int8}                     & \texttt{system.sbyte}    & 8-bit signed integer. \\
  \texttt{uint8}                    & \texttt{system.byte}     & 8-bit unsigned integer. \\
  \texttt{int16}                    & \texttt{system.int16}    & 16-bit signed integer. \\
  \texttt{uint16}                   & \texttt{system.uint16}   & 16-bit unsigned integer. \\
  \texttt{int32}                    & \texttt{system.int32}    & 32-bit signed integer. \\
  \texttt{uint32}                   & \texttt{system.uint32}   & 32-bit unsigned integer. \\
  \texttt{int64}                    & \texttt{system.int64}    & 64-bit signed integer. \\
  \texttt{uint64}                   & \texttt{system.uint64}   & 64-bit unsigned integer. \\
  \texttt{float32}, \texttt{single} & \texttt{system.single}   & 32-bit IEEE 754 float. \\
  \texttt{float64}, \texttt{double} & \texttt{system.double}   & 64-bit IEEE 754 float. \\
  \texttt{char}                     & \texttt{system.char}     & Unicode scalar (UTF-16 unit). \\
  \texttt{string}                   & \texttt{system.string}   & Immutable Unicode character sequence. \\
  \texttt{object}                   & \texttt{system.object}   & Top type; any value. \\
  \texttt{decimal}                  & \texttt{system.decimal}  & High-precision fixed-point decimal. \\
  \texttt{datetime}                 & \texttt{system.datetime} & Date and time instant. \\
  \texttt{timespan}                 & \texttt{system.timespan} & Duration / time interval. \\
  \texttt{guid}                     & \texttt{system.guid}     & 128-bit globally unique identifier. \\
  \bottomrule
\end{longtable}

\subsection{User-Defined Types}

A type defined in the module is either a \textbf{class}, \textbf{interface},
\textbf{struct}, or \textbf{enum}. At runtime, instances of user-defined
types are \emph{managed objects}: heap-allocated records that carry:

\begin{itemize}[leftmargin=1.8em]
  \item A \textbf{type name} identifying the declared type.
  \item A \textbf{field store}: a mapping from field name to value.
  \item An optional \textbf{host backing object} supplied by the host runtime
    (see Section~\ref{sec:hostinterop}).
  \item Zero or more \textbf{attribute instances} attached at load time.
\end{itemize}

\subsection{Inheritance}

A class may declare at most one \emph{base class}. A class or struct may
declare any number of \emph{implemented interfaces}. Runtimes are 
required to perform inheritance-aware type checks in V3
(see Section~\ref{sec:limits}).

\subsection{Null}

\texttt{null} is a valid value for any non-primitive type. The
\opcode{ldnull} instruction places it on the stack.

\subsection{Value Coercion Semantics}
\label{sec:coerce}

Because the evaluation stack carries untyped values, the runtime applies
the following coercion rules at instruction boundaries:

\begin{longtable}{lp{9.5cm}}
  \toprule
  \textbf{Target type} & \textbf{Rule} \\
  \midrule
  \endhead
  \texttt{bool}      &
    \texttt{null}, integer \texttt{0}, and empty string \texttt{""} are
    \texttt{false}; all other values are \texttt{true}. \\
  integer type       &
    Floating-point values are truncated toward zero. Strings are parsed
    as base-10 integers using the invariant locale. Out-of-range results
    are implementation-defined but must not raise an unhandled exception. \\
  floating-point     &
    Integer values are widened exactly. Strings are parsed as decimal
    floating-point using the invariant locale. \\
  \texttt{string}    &
    Any non-null value is converted to its canonical string representation.
    \texttt{null} converts to \texttt{null}. \\
  \texttt{object}    &
    Identity; any value is valid. \\
  \bottomrule
\end{longtable}

Arithmetic instructions select \emph{float mode} (operands widened to
64-bit float) when either operand is a floating-point value or a string;
otherwise \emph{integer mode} (64-bit signed) is used.

% ══════════════════════════════════════════════════════════════════════════════
\section{Instruction Set Reference}
\label{sec:isa}
% ══════════════════════════════════════════════════════════════════════════════

ObjectIR uses a \textbf{stack machine} execution model. Each method
invocation creates a new \emph{call frame} that owns:

\begin{itemize}[leftmargin=1.8em, noitemsep]
  \item An \textbf{evaluation stack} of untyped values.
  \item A named \textbf{argument table} (index \texttt{0} / name
    \texttt{this} is the receiver for instance methods).
  \item A named \textbf{local variable table}.
  \item A \textbf{pending exception} slot.
  \item A \textbf{program counter} (index into the instruction list).
\end{itemize}

Operands are described using the JSON operand format of
Section~\ref{sec:schema}.

\subsection{Stack Manipulation}

\begin{longtable}{lp{10.5cm}}
  \toprule
  \textbf{Opcode} & \textbf{Description} \\
  \midrule
  \endhead
  \opcode{nop}    & No-operation; advances the program counter only. \\
  \opcode{dup}    & Duplicate the top value on the evaluation stack. \\
  \opcode{pop}    & Discard the top value. \\
  \opcode{ldnull} & Push \texttt{null} onto the stack. \\
  \bottomrule
\end{longtable}

\subsection{Load Constants}

\begin{longtable}{lp{10.5cm}}
  \toprule
  \textbf{Opcode} & \textbf{Description} \\
  \midrule
  \endhead
  \opcode{ldc.i4} & Push a 32-bit integer constant. Operand: bare integer. \\
  \opcode{ldc.i8} & Push a 64-bit integer constant. Operand: bare integer. \\
  \opcode{ldc.r4} & Push a 32-bit float constant. Operand: bare number. \\
  \opcode{ldc.r8} & Push a 64-bit float constant. Operand: bare number. \\
  \opcode{load.str}  &
    Push a string constant.
    Operand: \code{"<Value>"} \\
  \bottomrule
\end{longtable}

\subsection{Arguments and Local Variables}

\begin{longtable}{lp{10.5cm}}
  \toprule
  \textbf{Opcode} & \textbf{Description} \\
  \midrule
  \endhead
  \opcode{load.arg}  &
    Push a method argument.
    Operand: parameter name or 0-based index.
    Index \texttt{0} / name \texttt{this} is the receiver. \\
  \opcode{store.arg}  & Pop a value and store it in the named argument slot. \\
  \opcode{load.loc}  & Push a local variable by name. \\
  \opcode{store.loc}  & Pop a value and store it in the named local variable. \\
  \bottomrule
\end{longtable}

\subsection{Field Access}

\begin{longtable}{lp{10.5cm}}
  \toprule
  \textbf{Opcode} & \textbf{Description} \\
  \midrule
  \endhead
  \opcode{load.fld}  &
    Pop an object (or use the implicit receiver) and push the value of the
    named instance field. If the object has a host backing, the runtime
    \textit{should} read the corresponding property on the host object
    before consulting the field store. \\
  \opcode{store.fld}  &
    Pop a value, then pop an object; store the value in the named instance
    field (and on the host backing, if present). \\
  \opcode{load.sfld} &
    Push the value of a static field.
    Operand: \code{<Field name>}. \\
  \opcode{store.sfld} & Pop a value and store it in the named static field. \\
  \bottomrule
\end{longtable}

\subsection{Arithmetic}

\begin{longtable}{lp{10.5cm}}
  \toprule
  \textbf{Opcode} & \textbf{Description} \\
  \midrule
  \endhead
  \opcode{add} &
    Pop $b$, pop $a$, push $a + b$.
    In float mode, also performs string concatenation when either operand
    is a string. \\
  \opcode{sub} & Pop $b$, pop $a$, push $a - b$. \\
  \opcode{mul} & Pop $b$, pop $a$, push $a \times b$. \\
  \opcode{div} & Pop $b$, pop $a$, push $a \div b$. \\
  \opcode{rem} & Pop $b$, pop $a$, push $a \bmod b$. \\
  \opcode{neg} & Pop $v$, push $-v$. \\
  \bottomrule
\end{longtable}

Mode selection follows Section~\ref{sec:coerce}: if either operand is
floating-point (or a string for \opcode{add}), both are promoted to 64-bit
float; otherwise 64-bit signed integer arithmetic is used.

\subsection{Logical}

\begin{longtable}{lp{10.5cm}}
  \toprule
  \textbf{Opcode} & \textbf{Description} \\
  \midrule
  \endhead
  \opcode{not} & Pop $v$, push the boolean negation: \texttt{!ToBool($v$)}. \\
  \bottomrule
\end{longtable}

\subsection{Comparison}

Each opcode pops $b$ then $a$ and pushes a \texttt{bool}.

\begin{longtable}{ll}
  \toprule
  \textbf{Opcode} & \textbf{Condition} \\
  \midrule
  \endhead
  \opcode{comp.eq} & $a = b$ \\
  \opcode{comp.ne} & $a \neq b$ \\
  \opcode{comp.gt} & $a > b$ \\
  \opcode{comp.ge} & $a \geq b$ \\
  \opcode{comp.lt} & $a < b$ \\
  \opcode{comp.le} & $a \leq b$ \\
  \bottomrule
\end{longtable}

String comparisons use ordinal (byte-by-byte) ordering.
Boolean values are supported for \opcode{ceq} and \opcode{cne} only.
All other operand pairs are promoted to 64-bit float before comparison.

\subsection{Type Conversion and Testing}

\begin{longtable}{lp{10.5cm}}
  \toprule
  \textbf{Opcode}    & \textbf{Description} \\
  \midrule
  \endhead
  \opcode{conv}      &
    Pop a value and push it converted to the named target type according
    to the coercion rules in Section~\ref{sec:coerce}. \\
  \opcode{castclass} &
    Assert that the top-of-stack managed object has the given type name;
    raise a cast error if it does not.) \\
  \opcode{isinst}    &
    Pop a value; push \texttt{true} if it is a managed object with the
    given type name, otherwise \texttt{false}. \\
  \bottomrule
\end{longtable}

\begin{infobox}
\textbf{V3 Info:} \opcode{castclass} and \opcode{isinst} perform
\emph{Rough} type-name matching. Walking an inheritance chain is not required
of conforming implementations in this version but is strongly recommended.
\end{infobox}

\subsection{Object and Array Creation}

\begin{longtable}{lp{10.5cm}}
  \toprule
  \textbf{Opcode}  & \textbf{Description} \\
  \midrule
  \endhead
  \opcode{new.obj}  &
    Create a new managed object of the named type and push it. If the type
    is a registered host type, the runtime \textit{should} also construct
    the corresponding host backing object. \\
  \opcode{new.arr}  &
    Push a new, empty, dynamically-typed array. \\
    optionaly with a type arg to set the type of the array.
    '\opcode{newarr int}'
  \opcode{load.elem}  &
    Pop an index (0-based integer), pop an array;
    push \texttt{array[index]}. \\
  \opcode{store.elem}  &
    Pop a value, pop an index, pop an array;
    set \texttt{array[index] = value}. \\
  \bottomrule
\end{longtable}

\begin{notebox}
Arrays in ObjectIR v3 are Static typed and variable-length. Typed array variants
(e.g.\ \texttt{int[]}) are Available in section~\ref{sec:Types}.
\end{notebox}

\subsection{Method Calls}

\begin{longtable}{lp{10.5cm}}
  \toprule
  \textbf{Opcode}   & \textbf{Description} \\
  \midrule
  \endhead
  \opcode{call}     &
    Pop arguments left-to-right, then pop the receiver (for instance
    methods); invoke the target method.
    Operand: \code{"<Namespace>.*(<namespace>.*).<functionName>(<Args>) -> <return type>"}. \\
  \opcode{callvirt} &
    Same as \opcode{call}, but dispatch is virtual: the actual type of the
    receiver determines which override is executed. \\
  \opcode{ret}      &
    Return from the current frame. If a value is present on the evaluation
    stack it is transferred to the caller's stack. \\
  \bottomrule
\end{longtable}

\subsubsection{Call Resolution Order}
\label{sec:callres}

A conforming runtime must attempt the following sequence for every
\opcode{call} and \opcode{callvirt}:

\begin{enumerate}[leftmargin=2em]
  \item \textbf{Exact native/host lookup.} Match the full signature string
    \code{DeclaringType.Name(param0,param1,...)} against the registered
    host method table.
  \item \textbf{Simplified native/host lookup.} Retry with parameter type
    names reduced to their unqualified, lower-case short names
    (e.g.\ \code{System.String} $\to$ \code{string}).
  \item \textbf{Host reflection.} If the runtime supports reflective dispatch
    and it is enabled, invoke the method on the host backing object.
  \item \textbf{IR method dispatch.} Look up the method in the loaded
    module's type table.
\end{enumerate}

If none of the steps resolve the call, a \textbf{missing-method} error is
raised.

\subsection{Control Flow}

\begin{longtable}{lp{10.5cm}}
  \toprule
  \textbf{Opcode}    & \textbf{Description} \\
  \midrule
  \endhead
  \opcode{if}        &
    Evaluate the condition operand (Section~\ref{sec:conditions}).
    If truthy, execute \code{thenBlock};
    otherwise execute the optional \code{elseBlock}. \\
  \opcode{while}     &
    Repeatedly evaluate the condition and execute \code{body}.
    A \opcode{break} inside \code{body} exits the loop;
    a \opcode{continue} skips to the next condition evaluation. \\
  \opcode{break}     & Exit the nearest enclosing \opcode{while}. \\
  \opcode{continue}  & Skip to the next iteration of the nearest \opcode{while}. \\
  \opcode{try}       &
    Execute \code{tryBlock}. If an exception is raised, test each
    \code{catchBlock} in order (Section~\ref{sec:exceptions}).
    Execute the optional \code{finallyBlock} in all cases. \\
  \opcode{throw}     &
    Pop a value and raise it as a pending exception on the current frame. \\
  \bottomrule
\end{longtable}

\subsubsection{Condition Kinds}
\label{sec:conditions}

A condition embedded in an \opcode{if} or \opcode{while} operand has a
\texttt{kind} field that determines how it is evaluated:

\begin{longtable}{lp{10cm}}
  \toprule
  \textbf{Kind}       & \textbf{Evaluation} \\
  \midrule
  \endhead
  \code{stack}      &
    Pop the top-of-stack value and coerce it to \texttt{bool}
    (Section~\ref{sec:coerce}).
    An empty stack is treated as \texttt{false}. \\
  \code{binary}     &
    Pop $r$, pop $l$; apply the named comparison opcode
    (\opcode{ceq}, \opcode{cgt}, \ldots) and use the result. \\
  \code{expression} &
    Execute a single embedded instruction and pop the result. \\
  \code{block}      &
    Execute a sequence of instructions and pop the final result. \\
  \bottomrule
\end{longtable}

% ══════════════════════════════════════════════════════════════════════════════
\section{Exception Model}
\label{sec:exceptions}
% ══════════════════════════════════════════════════════════════════════════════

\subsection{Raising Exceptions}

The \opcode{throw} instruction accepts any value: a managed object, a host
exception type, or a plain string. The value is stored in the
\emph{pending exception} slot of the current frame and instruction execution
in that frame halts.

\subsection{Structured Exception Handling}

The \opcode{try} instruction provides first-class structured exception
handling. Its operand defines up to three regions:

\begin{itemize}[leftmargin=1.8em]
  \item \code{tryBlock} --- the guarded instructions.
  \item \code{catchBlocks} --- an ordered list of catch clauses, each with
    an optional \code{type} field and a \code{body}.
  \item \code{finallyBlock} --- instructions executed unconditionally after
    the try/catch, regardless of outcome.
\end{itemize}

Catch clauses are matched in declaration order against the pending exception:

\begin{enumerate}[leftmargin=2em]
  \item \textbf{Catch-all} (no \code{type} field) --- matches any value.
  \item \textbf{Type name match} --- the pending exception is a managed object
    whose type name equals the catch clause's \code{type}.
  \item \textbf{Host type match} --- the host runtime may additionally apply
    assignability checks against registered host types.
\end{enumerate}

\subsection{Exception Propagation}

If no catch clause in the current frame matches, the pending exception
propagates to the caller's frame. This continues until either a matching
clause is found or the call stack is exhausted. An unhandled exception at
the top of the stack is delivered to the host via the exception hook.

\subsection{Control-Flow Exception Tokens}

The \opcode{break} and \opcode{continue} instructions are implemented as
\emph{structured control-flow tokens}. They must not be catchable by user
\opcode{try}/\opcode{catch} blocks and must not propagate outside the
enclosing \opcode{while} instruction.

% ══════════════════════════════════════════════════════════════════════════════
\section{FOB/IR v3 Binary Format}
\label{sec:fob}
% ══════════════════════════════════════════════════════════════════════════════

The FOB/IR v3 format is a compact binary encoding of the module schema
defined in Section~\ref{sec:schema}. It is intended for distribution and
fast loading; the semantics are identical to the JSON and TextIR forms.

\subsection{File Layout}

\begin{longtable}{lp{9cm}}
  \toprule
  \textbf{Section} & \textbf{Contents} \\
  \midrule
  \endhead
  \textbf{Header}        &
    4-byte magic \texttt{0x464F4249} (\texttt{FOBI}), 1-byte format
    version (\texttt{0x03}), module name (length-prefixed UTF-8), version
    triple (three \texttt{uint16}). \\
  \textbf{String pool}   &
    A length-prefixed table of interned UTF-8 strings referenced by index
    throughout the rest of the file. \\
  \textbf{Type table}    &
    Sequential type records: kind byte, name index, namespace index,
    base type index, interface index list, field count + field records,
    method count + method records. \\
  \textbf{Method bodies} &
    Per-method: parameter count + parameter records, local count + local
    records, instruction count + instruction records. \\
  \bottomrule
\end{longtable}

\subsection{Binary Opcode Map}

Each instruction is stored as a 1-byte opcode followed by a variable-length
operand.

\begin{longtable}{llll}
  \toprule
  \textbf{Byte} & \textbf{Mnemonic} & \textbf{Byte} & \textbf{Mnemonic} \\
  \midrule
  \endhead
  \texttt{0x00} & \texttt{nop}       & \texttt{0x10} & \texttt{load.sfld}   \\
  \texttt{0x01} & \texttt{ldc}       & \texttt{0x11} & \texttt{store.sfld}   \\
  \texttt{0x02} & \texttt{load.str}     & \texttt{0x12} & \texttt{new.obj}   \\
  \texttt{0x03} & \texttt{load.arg}     & \texttt{0x13} & \texttt{new.arr}   \\
  \texttt{0x04} & \texttt{store.arg}     & \texttt{0x14} & \texttt{load.elem}   \\
  \texttt{0x05} & \texttt{load.loc}     & \texttt{0x15} & \texttt{store.elem}   \\
  \texttt{0x06} & \texttt{store.loc}     & \texttt{0x16} & \texttt{call}     \\
  \texttt{0x07} & \texttt{add}       & \texttt{0x17} & \texttt{call.virt} \\
  \texttt{0x08} & \texttt{sub}       & \texttt{0x18} & \texttt{return}      \\
  \texttt{0x09} & \texttt{mul}       & \texttt{0x19} & \texttt{if}       \\
  \texttt{0x0A} & \texttt{div}       & \texttt{0x1A} & \texttt{while}    \\
  \texttt{0x0B} & \texttt{rem}       & \texttt{0x1B} & \texttt{break}    \\
  \texttt{0x0C} & \texttt{neg}       & \texttt{0x1C} & \texttt{continue} \\
  \texttt{0x0D} & \texttt{comp.eq}   & \texttt{0x1D} & \texttt{try}      \\
  \texttt{0x0E} & \texttt{comp.ne}   & \texttt{0x1E} & \texttt{throw}    \\
  \texttt{0x0F} & \texttt{load.fld}  & \texttt{0x1F} & \texttt{convert}     \\
  \texttt{0x20} & \texttt{castclass} & \texttt{0x21} & \texttt{isinst}   \\
  \texttt{0x22} & \texttt{dup}       & \texttt{0x23} & \texttt{pop}      \\
  \texttt{0x24} & \texttt{load.null}    & \texttt{0x25} & \texttt{not}      \\
  \texttt{0x26} & \texttt{comp.gt}       & \texttt{0x27} & \texttt{cge}      \\
  \texttt{0x28} & \texttt{comp.lt}       & \texttt{0x29} & \texttt{cle}      \\
  \texttt{0x2A} & \texttt{store.fld}     & \texttt{0x2B} & \texttt{ldc.i4}   \\
  \texttt{0x2C} & \texttt{ldc.i8}    & \texttt{0x2D} & \texttt{ldc.r4}   \\
  \texttt{0x2E} & \texttt{ldc.r8}    &               &                   \\
  \bottomrule
\end{longtable}

% ══════════════════════════════════════════════════════════════════════════════
\section{Host Interoperability}
\label{sec:hostinterop}
% ══════════════════════════════════════════════════════════════════════════════

This section defines requirements for runtimes that allow host code to
expose native types and methods to IR programs. All features in this section
are \textbf{optional}; a minimal runtime that only executes pure IR code
need not implement them.

\subsection{Host Type Registration}

A runtime \textit{may} allow the host to register \emph{external types}
under IR type names. Once registered:

\begin{itemize}[leftmargin=1.8em]
  \item \opcode{newobj} for that type name constructs a host instance and
    attaches it as the backing of a MethodNode for use within the runtime.
  \item \opcode{ldfld} / \opcode{stfld} consult the host object's property
    store before the managed object's field dictionary.
  \item \opcode{ldsfld} / \opcode{stsfld} may access static members on the
    host type.
\end{itemize}

\subsection{Native Method Bindings}

A runtime \textit{may} allow the host to register \emph{native method
bindings}: callbacks keyed by a call signature string of the form:

\[
\texttt{DeclaringType.MethodName(param0Type,param1Type,...)}
\]

When the runtime resolves a \opcode{call} or \opcode{callvirt}, it checks
the native binding table before performing IR method lookup (steps 1--3 of
Section~\ref{sec:callres}).

A native binding callback receives:
\begin{itemize}[leftmargin=1.8em, noitemsep]
  \item The receiver object (or a null equivalent for static calls).
  \item A positional list of argument values in declaration order.
\end{itemize}

It returns a single result value (or a null equivalent for void methods).

\subsection{Observability Hooks}

A runtime \textit{should} provide the following hooks for diagnostics and
tooling:

\begin{itemize}[leftmargin=1.8em]
  \item \textbf{Step hook} --- invoked before each instruction with the
    current frame and instruction as arguments. Useful for debuggers,
    profilers, and instruction-level traces.
  \item \textbf{Exception hook} --- invoked when an exception escapes the
    top of the call stack.
\end{itemize}

% ══════════════════════════════════════════════════════════════════════════════
\section{Runtime Conformance Requirements}
\label{sec:conformance}
% ══════════════════════════════════════════════════════════════════════════════

A conforming ObjectIR V3 runtime \textbf{must}:

\begin{enumerate}[leftmargin=2em]
  \item Accept at least one of the three module formats
    (TextIR, BSON, FOB/IR v3) as input.
  \item Maintain a call stack of frames with independent evaluation stacks.
  \item Implement all opcodes defined in Section~\ref{sec:isa}.
  \item Apply the coercion rules of Section~\ref{sec:coerce} at instruction
    boundaries.
  \item Implement the call resolution order of Section~\ref{sec:callres}.
  \item Implement structured exception handling as defined in
    Section~\ref{sec:exceptions}, including the requirement that
    \opcode{break}/\opcode{continue} tokens are not user-catchable.
  \item Expose an entry-point invocation mechanism (canonically targeting a
    static \texttt{Program.Main} with no parameters). see section see section \ref{sec:methodresolution}
  \item Expose a named method invocation mechanism allowing the host to call
    any method by qualified name with arbitrary arguments.
  \item Expose a named method existence check that inspects module metadata
    without executing code.
\end{enumerate}

A conforming runtime \textit{should} additionally:

\begin{enumerate}[leftmargin=2em, resume]
  \item Implement host type registration (Section~\ref{sec:hostinterop}).
  \item Implement native method binding (Section~\ref{sec:hostinterop}).
  \item Implement step and exception observability hooks.
  \item Pre-register the built-in \texttt{ObjectIR.Stdlib} module.
  for more information. see section \ref{sec:generics}
\end{enumerate}

% ══════════════════════════════════════════════════════════════════════════════
\section{Worked Examples}
% ══════════════════════════════════════════════════════════════════════════════

\subsection{Hello World}

\begin{lstlisting}[style=textir, caption={Hello World in TextIR}]
module HelloWorld version 1.0.0

class Program {
    static method Main() -> void {
        ldstr "Hello, ObjectIR!"
        call System.Console.WriteLine(string) -> void
        ret
    }
}
\end{lstlisting}

\subsection{Fibonacci (Recursive)}

\begin{lstlisting}[style=textir, caption={Recursive Fibonacci}]
module Fibonacci version 1.0.0

class Program {
    static method Fib(n: int32) -> int32 {
        ldarg n
        ldc.i4 2
        clt
        if (stack) 
        {
         	ldarg n
         	ret
        }

        ldarg n
        ldc.i4 1
        sub
        call Program.Fib(int32) -> int32

        ldarg n
        ldc.i4 2
        sub
        call Program.Fib(int32) -> int32

        add
        ret
    }

    static method Main() -> void {
        ldc.i4 10
        call Program.Fib(int32) -> int32
        call System.Console.WriteLine(int32) -> void
        ret
    }
}
\end{lstlisting}

\subsection{Instance Classes and Fields}

\begin{lstlisting}[style=textir, caption={Instance class with constructor}]
module ClassDemo version 1.0.0

class Vector2D {
    field x: float32
    field y: float32

    constructor(x: float32, y: float32) {
        ldarg this
        ldarg x
        stfld x
        ldarg this
        ldarg y
        stfld y
        ret
    }

    // returns sqrt(x*x + y*y)
    method Length() -> float32 {
        ldarg this
        ldfld x
        ldarg this
        ldfld x
        mul
        ldarg this
        ldfld y
        ldarg this
        ldfld y
        mul
        add
        call System.Math.Sqrt(float64) -> float64
        conv float32
        ret
    }
}

class Program {
    static method Main() -> void {
        ldc.r4 3.0
        ldc.r4 4.0
        newobj Vector2D
        callvirt Vector2D.Length() -> float32
        call System.Console.WriteLine(object) -> void
        ret
    }
}
\end{lstlisting}

\subsection{Exception Handling}

\begin{lstlisting}[style=textir, caption={try / catch / finally}]
module ExceptionDemo version 1.0.0

class MyError {
    field message: string
    constructor(message: string) {
        ldarg this
        ldarg message
        stfld message
        ret
    }
}

class Program {
    static method Main() -> void {
        try {
            newobj ExceptionDemo.MyError
            throw
            }
            catch (err: ExceptionDemo.MyError)
            {
            	ldstr "Caught!"
            	call IO.Println(string) -> void
            }
            finaly {
            	ldstr "Always runs"
            	call IO.Println(string) -> void
            	
            }
     
        ret
    }
}
\end{lstlisting}


\begin{lstlisting}[style=textir, caption={try / catch / finally}]
	module ExceptionDemo version 1.0.0
	
	class MyError {
		field message: string
		constructor(message: string) {
			ldarg this
			ldarg message
			stfld message
			ret
		}
	}
	
	class Program {
		static method Main() -> void {
			try {
				newobj ExceptionDemo.MyError
				throw
			}
			catch (err: ExceptionDemo.MyError)
			{
				ldstr "Caught!"
				call IO.Println(string) -> void
			}
			finaly {
				ldstr "Always runs"
				call IO.Println(string) -> void
				
			}
			
			ret
		}
	}
\end{lstlisting}

\section{Generics}
\label{sec:Types}
\label{sec:generics}
With OOP languages like CSharp or java, developers have access to what they call Generics
These generics allow for developers to make custom classes that don't rely on a set type to function.

unlike interfaces which define a concrete type that can be used to "hold" types. 
generics allow for developers to do whacky stuff.

\subsection{Generics preview}
ObjectIR V3 supports Generics via the generic attribute which exposes a changable type to the runtime.
this allow's developers to write programs in their host languages that preform some whacky stuff.

\begin{lstlisting}[style=textir, caption={generics example}]
module ExceptionDemo version 1.0.0

@Generic(T)
class List {
    field contents: T[]
    constructor(cts: T[]) {
        ldarg this
        ldarg cts
        stfld contents
        ret
    }
}

class Program {
    static method Main() -> void {
        try {
            	newobj ExceptionDemo.List<string>
            }
            catch (err: GenericConstructionError)
            {
            	ldstr "Caught!"
            	call IO.Println(string) -> void
            	call IO.Println(string) -> void // this works because when err is catched, it's on the stack still.
            	// we just then print it out
            }
            finaly {
            	ldstr "Always runs"
            	call IO.Println(string) -> void
            	
            }
     
        ret
    }
}
\end{lstlisting}

\section{Standard library}
\label{sec:stdlib}
ObjectIR now comes with a featured* standard library.

this standard library allows developers to have a proper API to target when developing tooling and such for ObjectIR.

you can check out the partnered implementation of ObjectIR.Stdlib
https://github.com/OpenStudioCorp/ObjectIR.Stdlib

*ObjectIR's standard library does not contain all functions required to 
self host infrastructure. this is constantly evolving.

% ══════════════════════════════════════════════════════════════════════════════
\section{Known Limitations (V2 -> V3)}
\label{sec:limits}
% ══════════════════════════════════════════════════════════════════════════════

The following features are explicitly \textbf{out of scope} for ObjectIR V2
But are being worked on for V3:

\begin{itemize}[leftmargin=1.8em]
  \item \textbf{Generics.} Generic type parameters and instantiation are
  available to be used, however the runtime must support Generics. 
  see section \ref{sec:generics}
    
  \item \textbf{Typed arrays.} All arrays are typed. Distinct
    \texttt{int[]}, \texttt{float[]}, etc. are available for use. 
    see section \ref{sec:arrays}
  
  \item \textbf{Inheritance-aware type tests.} \opcode{castclass} and
    \opcode{isinst} match by exact type name only.
    this has been superseeded, see section \ref{sec:generics}
    
  \item \textbf{Method overloading on IR-defined types.} Only one method
    with a given name is permitted per IR type in v2; overloading resolution
    applies only to host native bindings.
    this has been fixed, see section \ref{sec:methodresolution}

  \item \textbf{Module-level free functions.} The \texttt{functions} field does not exist on later versions.
  if your compier exports functions like this, they will NOT transfer over due to 
  the way the new ModuleNode system works in the latest version of ObjectIR.Core
  
  \item \textbf{Concurrent execution.} Thread safety is not a top priority for ObjectIR's runtime.
  Finite's lattice runtime supports* multithreading in a ObjectIR application.
\end{itemize}

% ══════════════════════════════════════════════════════════════════════════════
\appendix
% ══════════════════════════════════════════════════════════════════════════════


\section{Arrays}
\label{sec:arrays}
ObjectIR V3 now supports array's via the \opcode{ldelem} and \opcode{stelem} instructions.

creating array's are simple.
use \opcode{newarr}, optionaly followed by a type (eg \opcode{newarr int}) to make it a typed array instead of a dynamic-typed one.

store it into a local with \opcode{stloc}
then use \opcode{stelem} and \opcode{ldelem} to store items in arrays as follows

\begin{lstlisting}[style=textir, caption={array example}]
	module arraytest version 1.0.0
	class Program {
		static method Main() -> void {
			local arr: int[]
			
			newarr int // create the array
			
			ldc.i4 5  // store 5 into the array
			ldelem 0  // load the index
			
			ldc.i4 10 
			ldelem 1 
			
			ldc.i4 15 
			ldelem 2
			
			ldc.i4 20 
			ldelem 3 
			
			stloc arr // store the array
		}
	}
\end{lstlisting}


\section{Method Resolution}
\label{sec:methodresolution}

In ObjectIR, Method Resolution is primarily based on certain rules.

\begin{itemize}
	\item is the method externally defined?
	\item is the method defined with overloading?
	\end{itemize}
	and more.
	
	for methods that have been defined externaly, their BlockStatement will be null and instead a NativeMethod will exist in the MethodNode containing type parameters.
	
\section{Opcode Quick Reference}
\label{app:opcodes}

\begin{longtable}{lll}
  \toprule
  \textbf{Mnemonic} & \textbf{Category} & \textbf{Stack effect} \\
  \midrule
  \endhead
  \opcode{nop}       & Stack      & --- \\
  \opcode{dup}       & Stack      & $v \to v,\, v$ \\
  \opcode{pop}       & Stack      & $v \to$ \\
  \opcode{ldnull}    & Stack      & $\to \texttt{null}$ \\
  \opcode{ldc}       & Constants  & $\to c$ \\
  \opcode{ldc.i4}    & Constants  & $\to c_{\text{i32}}$ \\
  \opcode{ldc.i8}    & Constants  & $\to c_{\text{i64}}$ \\
  \opcode{ldc.r4}    & Constants  & $\to c_{\text{f32}}$ \\
  \opcode{ldc.r8}    & Constants  & $\to c_{\text{f64}}$ \\
  \opcode{ldstr}     & Constants  & $\to s$ \\
  \opcode{ldarg}     & Variables  & $\to a$ \\
  \opcode{starg}     & Variables  & $v \to$ \\
  \opcode{ldloc}     & Variables  & $\to l$ \\
  \opcode{stloc}     & Variables  & $v \to$ \\
  \opcode{ldfld}     & Fields     & $\mathit{obj} \to v$ \\
  \opcode{stfld}     & Fields     & $\mathit{obj}, v \to$ \\
  \opcode{ldsfld}    & Fields     & $\to v$ \\
  \opcode{stsfld}    & Fields     & $v \to$ \\
  \opcode{add}       & Arithmetic & $a, b \to a{+}b$ \\
  \opcode{sub}       & Arithmetic & $a, b \to a{-}b$ \\
  \opcode{mul}       & Arithmetic & $a, b \to a{\times}b$ \\
  \opcode{div}       & Arithmetic & $a, b \to a{\div}b$ \\
  \opcode{rem}       & Arithmetic & $a, b \to a \bmod b$ \\
  \opcode{neg}       & Arithmetic & $v \to {-}v$ \\
  \opcode{not}       & Logical    & $v \to \neg v$ \\
  \opcode{ceq}       & Comparison & $a, b \to a{=}b$ \\
  \opcode{cne}       & Comparison & $a, b \to a{\neq}b$ \\
  \opcode{cgt}       & Comparison & $a, b \to a{>}b$ \\
  \opcode{cge}       & Comparison & $a, b \to a{\geq}b$ \\
  \opcode{clt}       & Comparison & $a, b \to a{<}b$ \\
  \opcode{cle}       & Comparison & $a, b \to a{\leq}b$ \\
  \opcode{conv}      & Type       & $v \to v'$ \\
  \opcode{castclass} & Type       & $v \to v$ (or error) \\
  \opcode{isinst}    & Type       & $v \to \texttt{bool}$ \\
  \opcode{newobj}    & Objects    & $\to \mathit{obj}$ \\
  \opcode{newarr}    & Objects    & $\to \mathit{arr}$ \\
  \opcode{ldelem}    & Objects    & $\mathit{arr}, i \to v$ \\
  \opcode{stelem}    & Objects    & $\mathit{arr}, i, v \to$ \\
  \opcode{call}      & Calls      & $\mathit{args}\ldots \to [r]$ \\
  \opcode{callvirt}  & Calls      & $\mathit{args}\ldots, \mathit{obj} \to [r]$ \\
  \opcode{ret}       & Calls      & $[v] \to$ (returns to caller) \\
  \opcode{if}        & Control    & (condition-dependent) \\
  \opcode{while}     & Control    & (condition-dependent) \\
  \opcode{break}     & Control    & --- \\
  \opcode{continue}  & Control    & --- \\
  \opcode{try}       & Control    & --- \\
  \opcode{throw}     & Control    & $v \to$ (raises exception) \\
  \bottomrule
\end{longtable}

\end{document}

